\documentclass{article}
\usepackage{amsmath, amssymb, amsthm}
\usepackage{hyperref}
\usepackage{geometry}
\geometry{margin=1in}

\title{Erd\H{o}s Problem \#700}
\date{Last updated: 2025-08-31}
\author{Source: erdosproblems.com}

\begin{document}
\maketitle

\noindent\textbf{Status:} OPEN \quad \textbf{No prize}

\noindent\textbf{Tags:} number theory, binomial coefficients

\medskip
\noindent\textbf{Problem Statement:}

Let\[f(n)=\min_{1<k\leq n/2}\textrm{gcd}\left(n,\binom{n}{k}\right).\]{UL}
{LI}Characterise those composite $n$ such that $f(n)=n/P(n)$, where $P(n)$ is the largest prime dividing $n$.{/LI}
{LI}Are there infinitely many composite $n$ such that $f(n)>n^{1/2}$?{/LI}
{LI} Is it true that, for every composite $n$,\[f(n) \ll_A \frac{n}{(\log n)^A}\]for every $A>0$?{/LI}
{/UL}

\bigskip
\noindent\textbf{Remarks:}

A problem of Erd\H{o}s and Szekeres. It is easy to see that $f(n)\leq n/P(n)$ for composite $n$, since if $j=p^k$ where $p^k\mid n$ and $p^{k+1}\nmid n$ then $\textrm{gcd}\left(n,\binom{n}{j}\right)=n/p^k$. This implies\[f(n) \leq (1+o(1))\frac{n}{\log n}.\]It is known that $f(n)=n/P(n)$ when $n$ is the product of two primes. Another example is $n=30$.

For the second problem, it is easy to see that for any $n$ we have $f(n)\geq p(n)$, where $p(n)$ is the smallest prime dividing $n$, and hence there are infinitely many $n$ (those $=p^2)$ such that $f(n)\geq n^{1/2}$.

\bigskip
\noindent\small{Source: \url{https://www.erdosproblems.com/700}}
\end{document}
